\subsection{Dual of a locally convex metrizable space}

PROPOSITION 2. --- Let E be a locally convex metrizable space and F its strong dual. The space F is complete, semi-barrelled and satisfies the following condition :

(DB) There exists a sequence \((\mathbf{A}_n)_{n\in \mathbb{N}}\) of bounded subsets of \(\mathbf{F}\) such that every bounded subset of \(\mathbf{F}\) is contained in one of the \(\mathbf{A}_n\) .

The space E is bornological (III, p.~12, prop. 2), hence its strong dual is complete (III, p.~23, cor. 1).

Let \((\mathrm{V}_{n})_{n\in\mathbf{N}}\) be a decreasing sequence of neighbourhoods of 0 in E, such that every neighbourhood of 0 in E contains one of the \(V_{n}\) . Let \(A_{n}\) be the polar of \(V_{n}\) in F. Since E is bornological, every bounded subset of F is equicontinuous (III, p.~22, prop. 10), therefore contained in one of the \(A_{n}\) . In other words, the space F satisfies the condition (DB).

We now show that F is semi-barrelled. Let \((\mathrm{U}_{n})_{n\in\mathbf{N}}\) be a sequence of convex, balanced and closed neighbourhoods of 0 in F. We assume that the set \(U=\bigcap_{n}U_{n}\) absorbs every bounded subset of F. We shall prove that U is a neighbourhood of 0 in F. For this, we shall construct, by induction on the integer \(n\geqslant0\) , real numbers \(\lambda_{n}>0\) and convex balanced neighbourhoods \(W_{n}\) of 0 in F, whose which are closed for \(\sigma(\mathrm{F},\mathrm{E})\) , and satisfy the following relations

\[
\lambda_ {n} \mathrm{A} _ {n} \subset \frac {1}{2} \mathrm{U} \cap \left(\bigcap_ {0 \leqslant i <   n} \mathrm{W} _ {i}\right)\tag{1}
\]

\[
\bigcup_ {0 \leqslant i \leqslant n} \lambda_ {i} \mathrm{A} _ {i} \subset \mathrm{W} _ {n} \subset \mathrm{U} _ {n}.\tag{2}
\]

Suppose that the numbers \(\lambda_{i}\) and the sets \(\mathrm{W}_i\) have been constructed for \(0 \leqslant i < n\) . By the hypothesis, the set U absorbs the bounded subsets of F; moreover, for \(0 \leqslant i < n\) , \(\mathrm{W}_i\) is a neighbourhood of 0 in F, hence absorbs the bounded subsets of F. We can therefore find a number \(\lambda_{n} > 0\) satisfying (1). Let C denote the closed convex balanced envelope, for \(\sigma(\mathrm{F}, \mathrm{E})\) , of \(\bigcup_{0 \leqslant i \leqslant n} \lambda_i \mathrm{A}_i\) ; the set C is equicontinuous,

hence compact for \(\sigma(\mathrm{F},\mathrm{E})\) (III, p.~17, cor. 2). Since \(\mathrm{U}_n\) is a neighbourhood of 0 in F, there exists a bounded subset B of E such that \(\mathbf{B}^{\circ}\subset\frac{1}{2}\mathbf{U}_{n}\) . Put \(\mathrm{W}_n=\mathrm{C}+\mathrm{B}^{\circ}\) . Since \(\mathbf{B}^{\circ}\) is a neighbourhood of 0 in F, we see that \(\mathrm{W}_n\) is a convex and balanced neighbourhood of 0 in F. In addition, C is compact and \(\mathbf{B}^{\circ}\) closed for \(\sigma(\mathrm{F},\mathrm{E})\) ; by cor. 1 of GT, III, § 4, No.~1, \(\mathrm{W}_n\) is closed for \(\sigma(\mathrm{F},\mathrm{E})\) . Finally, we have \(\mathbf{C}\subset\frac{1}{2}\mathbf{U}\subset\frac{1}{2}\mathbf{U}_n\) and \(\mathbf{B}^{\circ}\subset\frac{1}{2}\mathbf{U}_n\) , hence \(\mathrm{W}_n\subset\mathrm{U}_n\) since \(\mathrm{U}_n\) is convex. We have thus established (2).

Put
\[
\mathrm{W} = \bigcap_ {n} \mathrm{W} _ {n}
\]

then
\[
\mathbf {W} \subset \mathbf {U}
\]

By (1) and (2), we have
\[
\lambda_ {i} \mathrm{A} _ {i} \subset \mathrm{W} _ {j} \quad \text {for all } i \text { and } j
\]

in N, and so \(\lambda_{i}A_{i}\subset W\) for all \(i\in N\) . In particular, W is absorbent, hence is a barrel for \(\sigma(\mathrm{F},\mathrm{E})\) . By remark 3 of IV, p.~4, W is a neighbourhood of 0 in F. A fortiori, U is a neighbourhood of 0 in F, and F is semi-barrelled.

The following corollary extends the Banach-Steinhaus theorem to the dual of a Fréchet space (cf.~III, p.~25, cor. 2).

COROLLARY. --- Let G be a Hausdorff locally convex space, and let \((u_{n})\) be a sequence of linear mappings from F into G, converging simply to a mapping u from F into G. Then u is continuous, and the sequence \((u_{n})\) converges to u uniformly on every pre-compact subset of F.

Since F is complete, the set of all \(u_{n}\) , which is bounded for the topology of simple convergence, is bounded in \(\mathcal{L}_{b}(F;G)\) (III, p.~27, cor. 1). Since the space F is semibarrelled (prop. 2), every countable and bounded subset of \(\mathcal{L}_{b}(F;G)\) is equicontinuous by prop. 1 of IV, p.~21. Therefore the set of the \(u_{n}\) is equicontinuous, and the corollary follows from III, p.~18, corollary.

